\usepackage[
  a4paper,
  inner=30mm,
  outer=25mm,
  bindingoffset=2mm,
  top=24mm,
  bottom=26mm,
  includeheadfoot,
  headheight=22pt,
  headsep=14pt,
  footskip=23pt
]{geometry}

\usepackage{emptypage}
\usepackage[english]{babel}
\usepackage{fontspec}
\usepackage{titlesec}
\usepackage{amsmath,amssymb,amsfonts,amsthm,mathtools}
\usepackage{bm}
\usepackage{booktabs,longtable,array,multirow}
\usepackage{enumitem}
\usepackage{xcolor}
\usepackage{tikz}
\usetikzlibrary{arrows.meta,positioning,calc,backgrounds}
\usepackage{tikz-cd}
\usepackage[most]{tcolorbox}
\usepackage{fancyhdr}

\definecolor{OxfordNavy}{RGB}{24, 48, 65}
\definecolor{DefinitionInk}{RGB}{31, 88, 83}
\definecolor{DeepBurgundy}{RGB}{114, 50, 57}
\definecolor{WarmGold}{RGB}{146, 112, 64}
\definecolor{SlateInk}{RGB}{55, 66, 73}
\definecolor{CodeSlate}{RGB}{48, 62, 78}
\definecolor{LinkBlue}{RGB}{18, 86, 174}
\definecolor{SoftIvory}{RGB}{250, 249, 246}
\definecolor{SoftSlateBg}{RGB}{245, 247, 250}
\definecolor{HairlineGray}{RGB}{185, 192, 200}

\usepackage{xurl}
\usepackage{csquotes}
\usepackage{zref-perpage}
\zmakeperpage{footnote}
\usepackage[
  backend=biber,
  style=numeric-comp,
  sorting=nyt,
  giveninits=true,
  maxbibnames=6,
  doi=true,
  isbn=false,
  url=true
]{biblatex}
\addbibresource{../references.bib}
\DeclareSourcemap{
  \maps[datatype=bibtex]{
    \map[overwrite]{
      \step[fieldsource=entrykey, match=\regexp{\AAlpoge2026\z}, final]
      \step[fieldset=title, fieldvalue={A note on the {Riemann} hypothesis and its generalization to {Dirichlet} \(L\)-functions}]
    }
    \map[overwrite]{
      \step[fieldsource=entrykey, match=\regexp{\ASIYZ2026\z}, final]
      \step[fieldset=title, fieldvalue={Why did {OpenAI} not prove the {Riemann} hypothesis first? {Evidence} and timeline (in {Chinese})}]
      \step[fieldset=howpublished, fieldvalue={Zhihu Column}]
    }
  }
}
\defbibheading{monographbib}[\bibname]{%
  \chapter*{#1}%
  \addcontentsline{toc}{chapter}{#1}%
  \markboth{#1}{#1}%
}
\renewcommand*{\bibfont}{\small}
\renewcommand*{\UrlFont}{\ttfamily\footnotesize}

\usepackage[
  colorlinks=true,
  hypertexnames=false,
  linktoc=all,
  linkcolor=OxfordNavy,
  citecolor=DeepBurgundy,
  urlcolor=LinkBlue,
  bookmarksnumbered=true,
  pdftitle={From Theta Lifts to the Proof of the Generalized Riemann Hypothesis},
  pdfauthor={Based on the lsernau/grh-lean formalization}
]{hyperref}
\usepackage{bookmark}

% Text fonts for the English edition.
\setmainfont{texgyrepagella}[
  Extension=.otf, UprightFont=*-regular, ItalicFont=*-italic,
  BoldFont=*-bold, BoldItalicFont=*-bolditalic]
\setsansfont{texgyreheros}[
  Extension=.otf, UprightFont=*-regular, ItalicFont=*-italic,
  BoldFont=*-bold, BoldItalicFont=*-bolditalic]
\setmonofont{Latin Modern Mono}

\linespread{1.20}
\setlength{\parskip}{0pt}
\setlength{\emergencystretch}{2em}
\allowdisplaybreaks[2]

% Running headers (preserve proper capitalization of math names like Seesaw, Hecke)
\pagestyle{fancy}
\fancyhf{}
\makeatletter
\renewcommand{\chaptermark}[1]{\markboth{\ifnum\value{secnumdepth}>-1\if@mainmatter\@chapapp\ \thechapter\quad\fi\fi #1}{}}
\makeatother
\renewcommand{\sectionmark}[1]{\markright{\thesection\quad #1}}
\fancyhead[LE]{\small\bfseries\color{SlateInk}\leftmark}
\fancyhead[RO]{\small\bfseries\color{SlateInk}\rightmark}
\fancyfoot[LE,RO]{\footnotesize\color{SlateInk}\thepage}
\renewcommand{\headrulewidth}{0.35pt}
\renewcommand{\headrule}{\hbox to\headwidth{\color{HairlineGray}\leaders\hrule height \headrulewidth\hfill}}

\fancypagestyle{plain}{%
  \fancyhf{}%
  \fancyfoot[C]{\footnotesize\color{SlateInk}\thepage}%
  \renewcommand{\headrulewidth}{0pt}%
}

\titleformat{\chapter}[display]
  {\raggedright\huge\bfseries\color{OxfordNavy}}
  {\LARGE\sffamily\bfseries\color{DefinitionInk}\chaptertitlename\ \thechapter}
  {0.65em}{}
\titlespacing*{\chapter}{0pt}{0.7em}{1.15em}

\titleformat{\section}
  {\large\bfseries\color{OxfordNavy}\raggedright}{\thesection}{1em}{}
\titlespacing*{\section}{0pt}{1.45em plus 0.3em minus 0.2em}{0.65em}

\titleformat{\subsection}
  {\normalsize\bfseries\color{SlateInk}\raggedright}{\thesubsection}{1em}{}
\titlespacing*{\subsection}{0pt}{1.05em plus 0.2em minus 0.1em}{0.4em}

\titleformat{\subsubsection}
  {\normalsize\bfseries\color{SlateInk}}{\thesubsubsection}{1em}{}

% Compact scholarly list spacing
\setlist[itemize]{topsep=0.4em, itemsep=0.22em, parsep=0pt, leftmargin=2em}
\setlist[enumerate]{topsep=0.4em, itemsep=0.22em, parsep=0pt, leftmargin=2.2em}

% Subtle dark-slate Lean identifier macros; web links inherit LinkBlue
\DeclareRobustCommand{\lean}[1]{\ifmmode\text{\upshape\ttfamily\small\color{CodeSlate}#1}\else{\upshape\ttfamily\small\color{CodeSlate}#1}\fi}
\DeclareRobustCommand{\leanlink}[2]{\href{#1}{{\upshape\ttfamily\small #2}}}
\DeclareRobustCommand{\leanfile}[1]{\ifmmode\text{\upshape\ttfamily\footnotesize\color{SlateInk}#1}\else{\upshape\ttfamily\footnotesize\color{SlateInk}#1}\fi}
\newcommand{\leansourcelink}[2]{%
  \href{https://github.com/lsernau/grh-lean/blob/72fead102f663d5c7a407964a13dc0cf9d977207/grh/#1\#L#2}{{\upshape\ttfamily\nolinkurl{#1}\#L#2}}%
}
\newcommand{\leansource}[2]{\leansourcelink{#1}{#2}}
\newcommand{\leanfn}[1]{\footnote{\normalfont\footnotesize\color{SlateInk}#1}}

% Refined footnote rule and spacing
\renewcommand{\footnoterule}{%
  \kern-3pt
  {\color{HairlineGray}\hrule width 0.34\textwidth height 0.4pt}%
  \kern2.6pt
}
\setlength{\footnotesep}{0.45em}

\newcommand{\chapref}[1]{\hyperref[#1]{Chapter~\ref*{#1}}}
\newcommand{\secref}[1]{\hyperref[#1]{Section~\ref*{#1}}}
\newcommand{\appref}[1]{\hyperref[#1]{Appendix~\ref*{#1}}}
\newcommand{\statementlink}[2]{\hyperref[#1]{\textcolor{LinkBlue}{#2}}}
\newcommand{\statementref}[2]{\statementlink{#2}{#1~\ref*{#2}}}
\pdfstringdefDisableCommands{%
  \def\lean#1{#1}%
  \def\leanlink#1#2{#2}%
  \def\leanfile#1{#1}%
  \def\RePart{Re }%
}

% Tate-Shafarevich symbol Sha via crisp TeX rules
\newcommand{\Sha}{\mathord{\mkern1.5mu\rule{0.5pt}{1.3ex}\rule{0.36em}{0.5pt}\rule{0.5pt}{1.3ex}\rule{0.36em}{0.5pt}\rule{0.5pt}{1.3ex}\mkern1.5mu}}

% Standard Number Theory & Automorphic Macros
\newcommand{\NN}{\mathbb{N}}
\newcommand{\ZZ}{\mathbb{Z}}
\newcommand{\QQ}{\mathbb{Q}}
\newcommand{\RR}{\mathbb{R}}
\newcommand{\CC}{\mathbb{C}}
\newcommand{\AAf}{\mathbb{A}_f}
\newcommand{\AAA}{\mathbb{A}}
\newcommand{\FF}{\mathbb{F}}
\newcommand{\OO}{\mathcal{O}}
\newcommand{\pp}{\mathfrak{p}}
\newcommand{\qq}{\mathfrak{q}}
\newcommand{\mm}{\mathfrak{m}}
\newcommand{\aaIdeal}{\mathfrak{a}}
\newcommand{\pMap}{\mathfrak{p}}
\newcommand{\zVol}{\mathfrak{z}}
\newcommand{\Gbar}{\overline{G}}

\DeclareMathOperator{\RePart}{Re}
\DeclareMathOperator{\ImPart}{Im}
\DeclareMathOperator{\Res}{Res}
\DeclareMathOperator{\ord}{ord}
\DeclareMathOperator{\Gal}{Gal}
\DeclareMathOperator{\Aut}{Aut}
\DeclareMathOperator{\End}{End}
\DeclareMathOperator{\Hom}{Hom}
\DeclareMathOperator{\GL}{GL}
\DeclareMathOperator{\SL}{SL}
\DeclareMathOperator{\Ugrp}{U}
\DeclareMathOperator{\Span}{span}
\DeclareMathOperator{\im}{im}
\DeclareMathOperator{\tr}{tr}
\DeclareMathOperator{\rk}{rk}
\DeclareMathOperator{\vol}{vol}
\DeclareMathOperator{\covol}{covol}
\DeclareMathOperator{\IndexLe}{IndexLe}
\DeclareMathOperator{\CovolumeLe}{CovolumeLe}
\DeclareMathOperator{\HeightLe}{HeightLe}
\DeclareMathOperator{\ColengthDvd}{ColengthDvd}
\DeclareMathOperator{\GaugeLe}{GaugeLe}
\DeclareMathOperator{\DenominatorLe}{DenominatorLe}
\DeclareMathOperator{\Frob}{Frob}
\DeclareMathOperator{\Spec}{Spec}
\DeclareMathOperator{\Fil}{Fil}
\DeclareMathOperator{\Sym}{Sym}
\DeclareMathOperator{\diag}{diag}
\DeclareMathOperator{\HS}{HS}

% Airy, borderless-except-left-accent Chapter Abstract (Cambridge / EMS Press style)
\newtcolorbox{chapterguide}[1][Chapter overview]{
  enhanced,
  breakable,
  frame hidden,
  colback=white,
  borderline west={1pt}{0pt}{HairlineGray},
  sharp corners,
  left=9pt,right=2pt,top=4pt,bottom=4pt,
  before skip=0.2em,
  after skip=1.3em,
  fontupper=\small\color{SlateInk},
  title={\small\bfseries\color{OxfordNavy}#1},
  colbacktitle=white,
  boxed title style={empty},
  attach title to upper={\quad}
}

% Subtle left-accent block ONLY for Landmark Main Theorems (inequality (*), public theorems)
\newtcolorbox{landmarkbox}{
  enhanced,
  breakable,
  frame hidden,
  colback=SoftIvory,
  borderline west={2pt}{0pt}{DeepBurgundy},
  sharp corners,
  left=10pt,right=10pt,top=3pt,bottom=3pt,
  before skip=0.7em,
  after skip=0.7em
}

% Clean unboxed structural remark block
\newenvironment{keyinsight}[1][Structural remark]{%
  \par\medskip\noindent{\bfseries\color{OxfordNavy}#1.}\quad\ignorespaces
}{%
  \par\medskip
}

% Hook to attach primary Lean source links cleanly as theorem-header footnotes
\makeatletter
\newcommand{\@pendingthmsource}{}
\newcommand{\thmsource}[1]{\gdef\@pendingthmsource{#1}}
\newcommand{\@emitthmsourcemark}{%
  \ifx\@pendingthmsource\@empty\else
    \kern0.08em{\normalfont\footnotemark}%
  \fi
}
\newcommand{\emitthmsource}{%
  \ifx\@pendingthmsource\@empty\else
    \leavevmode
    \footnotetext{\normalfont\footnotesize\color{SlateInk}\@pendingthmsource}%
    \gdef\@pendingthmsource{}%
    \ignorespaces
  \fi
}

\newtheoremstyle{thmstyle}
  {0.75em}{0.75em}
  {\normalfont\itshape}
  {0pt}{\bfseries\color{OxfordNavy}}{.\@emitthmsourcemark}{0.5em}
  {\thmname{#1}\thmnumber{ #2}\thmnote{ {\normalfont\color{SlateInk}(#3)}}}

\newtheoremstyle{mainstyle}
  {0.6em}{0.6em}
  {\normalfont\itshape}
  {0pt}{\bfseries\color{DeepBurgundy}}{.\@emitthmsourcemark}{0.5em}
  {\thmname{#1}\thmnumber{ #2}\thmnote{ {\normalfont\color{SlateInk}(#3)}}}

\newtheoremstyle{defstyle}
  {0.7em}{0.7em}
  {\normalfont}
  {0pt}{\bfseries\color{DefinitionInk}}{.\@emitthmsourcemark}{0.5em}
  {\thmname{#1}\thmnumber{ #2}\thmnote{ {\normalfont\color{SlateInk}(#3)}}}

\newtheoremstyle{lemstyle}
  {0.7em}{0.7em}
  {\normalfont}
  {0pt}{\bfseries\color{SlateInk}}{.\@emitthmsourcemark}{0.5em}
  {\thmname{#1}\thmnumber{ #2}\thmnote{ {\normalfont(#3)}}}

\newtheoremstyle{remstyle}
  {0.6em}{0.6em}
  {\normalfont}
  {0pt}{\bfseries\color{SlateInk}}{.\@emitthmsourcemark}{0.5em}
  {\thmname{#1}\thmnumber{ #2}\thmnote{(#3)}}
\makeatother

\theoremstyle{mainstyle}
\newtheorem{maintheorem}{Main Theorem}[chapter]

\theoremstyle{thmstyle}
\newtheorem{theorem}{Theorem}[chapter]
\newtheorem{proposition}[theorem]{Proposition}
\newtheorem{corollary}[theorem]{Corollary}

\theoremstyle{lemstyle}
\newtheorem{lemma}[theorem]{Lemma}

\theoremstyle{defstyle}
\newtheorem{definition}[theorem]{Definition}
\newtheorem{hypothesis}[theorem]{Hypothesis}

\theoremstyle{remstyle}
\newtheorem{remark}[theorem]{Remark}
\newtheorem{example}[theorem]{Example}

% Proof header in bold upright type
\makeatletter
\renewenvironment{proof}[1][\proofname]{\par
  \pushQED{\qed}%
  \normalfont \topsep6\p@\@plus6\p@\relax
  \trivlist
  \item[\hskip\labelsep\bfseries\color{OxfordNavy}#1\@addpunct{.}]\ignorespaces
}{%
  \popQED\endtrivlist\@endpefalse
}
\makeatother
\renewcommand{\proofname}{Proof}
